\uuid{EMpt}
\titre{Produit}
\auteur{bodin ; hanani ; mzari}
\organisation{exo7}
\licence{CC BY-NC-SA 4.0}
\datecreate{2020-01-20}
\niveau{L1}
\module{Algèbre}
\chapitre{Polynôme, fraction rationnelle}
\sousChapitre{Définition, degré, produit}
\difficulte{1}
\theme{Polynômes, produit}
\exo7qcmid{lille-1.polynomes.002}

\contenu{
\begin{question}
\qclassification{105.05}
\qtags{motcle=produit}
Soit $P(X) = X^3-3X^2+2$ et $Q(X) = X^3-X+1$.
Quelles sont les assertions vraies ?
\begin{answers}
    \bad{Le polynôme $P(X) \times Q(X)$ est de degré $9$.}

    \bad{Le coefficient du monôme $X^2$ dans le produit $P(X) \times Q(X)$ est $3$.}

    \good{Le polynôme $P(X) + Q(X)$ est de degré $3$.}

    \bad{Le polynôme $P(X) - Q(X)$ est de degré $3$.}
\end{answers}
\begin{explanations}
$P(X)\times Q(X) = X^6 - 3 X^5 - X^4 + 6 X^3 - 3 X^2 - 2 X + 2$,
$P(X) + Q(X) = 2 X^3 - 3 X^2 - X + 3$,
$P(X) - Q(X) = -3 X^2 + X + 1$.
\end{explanations}
\end{question}
}
